\section*{分部积分法}

1. 求下列不定积分。

(1) \(\int x\cos x\,dx\)。

解：分部积分，令 \(u=x\)，\(dv=\cos x\,dx\)，则 \(du=dx\)，\(v=\sin x\)。
\[ \int x\cos x\,dx=x\sin x-\int\sin x\,dx=x\sin x+\cos x+C. \]

(2) \(\int\ln x\,dx\)。

解：分部积分，令 \(u=\ln x\)，\(dv=dx\)，得
\[ \int\ln x\,dx=x\ln x-\int x\cdot\frac{1}{x}dx=x\ln x-x+C. \]

(3) \(\int x^{2}e^{x}\,dx\)。

解：分部两次，最终得
\[ \int x^{2}e^{x}\,dx=e^{x}(x^{2}-2x+2)+C. \]

(4) \(\int\arcsin x\,dx\)。

解：分部积分，令 \(u=\arcsin x\)，\(dv=dx\)，则
\[ \int\arcsin x\,dx=x\arcsin x-\int\frac{x}{\sqrt{1-x^{2}}}dx=x\arcsin x+\sqrt{1-x^{2}}+C. \]

(5) \(\int\frac{\ln(\ln x)}{x}\,dx\)。

解：令 \(u=\ln x\)，则 \(dx=x\,du\)，得
\[ \int\frac{\ln(\ln x)}{x}\,dx=\int\ln u\,du=u\ln u-u+C=\ln x\cdot\ln(\ln x)-\ln x+C. \]

(6) \(\int e^{2x}\cos x\,dx\)。

解：设 \(I=\int e^{2x}\cos x\,dx\)，按常规两次分部并解方程组，得
\[ I=\frac{1}{5}e^{2x}(\sin x+2\cos x)+C. \]

(7) \(\int x\sin x\cos x\,dx\)。

解：利用 \(\sin x\cos x=\frac{1}{2}\sin 2x\)，再分部积分，得
\[ \int x\sin x\cos x\,dx=\frac{1}{2}\int x\sin 2x\,dx=-\frac{1}{4}x\cos 2x+\frac{1}{8}\sin 2x+C. \]

(8) \(\int xf''(x)\,dx\)。

解：由分部积分法，设
\[ u=x,\qquad dv=f''(x)\,dx, \]
则
\[ du=dx,\qquad v=f'(x). \]
根据分部积分公式得
\[ \int xf''(x)\,dx=xf'(x)-\int f'(x)\,dx. \]
继续计算右边的积分：
\[ \int f'(x)\,dx=f(x)+C_1. \]
因此
\[ \int xf''(x)\,dx=xf'(x)-f(x)+C. \]

(9) \(\int x\sin^{2}x\,dx\)。

解：用 \(\sin^{2}x=\frac{1-\cos 2x}{2}\) 并分部积分，得
\[ \int x\sin^{2}x\,dx=\frac{1}{4}x^{2}-\frac{1}{4}x\sin 2x-\frac{1}{8}\cos 2x+C. \]

(10) \(\int x(\arctan x)^{2}\,dx\)。

解：分部积分（令 \(u=(\arctan x)^{2}\)，\(dv=x\,dx\)）；化简后得
\[ \int x(\arctan x)^{2}\,dx=\frac{1}{2}x^{2}(\arctan x)^{2}+\frac{1}{2}(\arctan x)^{2}-x\arctan x+\frac{1}{2}\ln(1+x^{2})+C. \]

(11) \(\int\ln\left(x+\sqrt{1+x^{2}}\right)dx\)。

解：分部积分并代换，得
\[ \int\ln\left(x+\sqrt{1+x^{2}}\right)dx=x\ln\left(x+\sqrt{1+x^{2}}\right)-\sqrt{1+x^{2}}+C. \]

(12) \(\int\frac{x\cos x}{\sin^{3}x}\,dx\)。

解：分部积分，令 \(u=x\)，\(dv=\frac{\cos x}{\sin^{3}x}dx\)。注意
\[ v=\int\frac{\cos x}{\sin^{3}x}dx=-\frac{1}{2\sin^{2}x}. \]
因此
\[ \int\frac{x\cos x}{\sin^{3}x}\,dx=-\frac{x}{2\sin^{2}x}-\frac{1}{2}\cot x+C. \]

(13) \(\int\sec^{5}x\,dx\)。

解：使用递推与已知 \(\int\sec^{3}x\,dx\) 的结果，得
\[ \int\sec^{5}x\,dx=\frac{1}{4}\sec^{3}x\tan x+\frac{3}{8}\sec x\tan x+\frac{3}{8}\ln|\sec x+\tan x|+C. \]

(14) \(\int\frac{x^{2}\arctan x}{1+x^{2}}\,dx\)。

解：注意到
\[ \frac{x^{2}}{1+x^{2}}=1-\frac{1}{1+x^{2}}, \]
因此
\[ \int\frac{x^{2}\arctan x}{1+x^{2}}\,dx=\int\arctan x\,dx-\int\frac{\arctan x}{1+x^{2}}\,dx. \]
先计算第一个积分（分部积分）：令 \(u=\arctan x\)，\(dv=dx\)，则 \(du=\frac{dx}{1+x^{2}}\)，\(v=x\)，故
\[ \int\arctan x\,dx=x\arctan x-\int\frac{x}{1+x^{2}}dx=x\arctan x-\frac{1}{2}\ln(1+x^{2})+C_1. \]
再计算第二个积分（换元）：令 \(t=\arctan x\)，则 \(dt=\frac{dx}{1+x^{2}}\)，于是
\[ \int\frac{\arctan x}{1+x^{2}}\,dx=\int t\,dt=\frac{t^{2}}{2}+C_2=\frac{(\arctan x)^{2}}{2}+C_2. \]
合并两部分并吸收常数，得最终结果：
\[ \int\frac{x^{2}\arctan x}{1+x^{2}}\,dx=x\arctan x-\frac{1}{2}\ln(1+x^{2})-\frac{1}{2}(\arctan x)^{2}+C. \]

2. 对于正整数 \(n\ge 2\)，建立 \(I_n=\int\sin^{n}x\,dx\) 的递推公式。

解：按分部积分法，令 \(u=\sin^{n-1}x\)，\(dv=\sin x\,dx\)，可得
\[ I_n=-\sin^{n-1}x\cos x+(n-1)I_{n-2}-(n-1)I_n. \]
移项后得到递推关系：
\[ I_n=-\frac{1}{n}\sin^{n-1}x\cos x+\frac{n-1}{n}I_{n-2}+C. \]
并给出初值 \(I_0=x+C\)，\(I_1=-\cos x+C\)。
